\honest{To Spread Science, Keep It Secret}{Eliezer Yudkowsky}{2008}

Sometimes I wonder if the Pythagoreans had the right idea. \nb{The Pythagoreans are said to have kept their teachings secret. The \textsc{Stanford Encyclopedia of Philosophy} finds this attested for some doctrines, and traces the idea that all of them were secret to later tradition.}

I have written that science is public by nature. But imagine another future, in which science is locked away and guarded by robed cults, so that more people will actually study it. Right now people can study science, and don't.

The reason is ``Scarcity'': what seems hard to get is valued more, ``especially'' information. \nb{The studies of information in ``Scarcity,'' posted the day before, show forbidden or exclusive information gaining influence. None measures how hard people try to obtain it.} People assume, I think, that if information is freely available it must not be important, so they join cults that keep their Great Truths secret. I give no evidence for this.

If science were kept in vaults, people would be desperate to get in, and rival cults would have to explain why their masters cannot match the Eighth Level Physicists. You will object that science already has fearsome initiations and is hard to learn, so why does that not count? My answer: ``Because the public thinks that science is freely available, that's why.''

If the secret of natural selection were given only after a ceremony with torches, robes and an ox, you would be satisfied with it, and laugh at creationists. \nb{Aronson and Mills (1959) found that a severe initiation made a group seem more attractive to those who went through it. The finding supports the post.} It might be more fun, too, especially if initiates had to put the evidence together themselves.

I am not seriously proposing this, ``At least, not at the moment.'' I only want you to imagine the robes and masks, so that you do not undervalue science for being unprotected, and to trust no group that wears them ``unless they also show you the data.''

People seem to have holes in their minds for secret knowledge, and if true beliefs do not fill them, false ones will. Last, I admit that science is not really free: the courses and textbooks are expensive. \nb{Cost, like the difficulty granted earlier, is a plain reason why few people study science. The post gives no evidence that favors its own explanation over either.} Yet the public thinks anyone may know it. It should be the other way around.
